% =============================================================
% II. SOURCES AND METHOD
% =============================================================
\section{Sources and method}
\label{sec:method}

\subsection{The production approach and the reliability of Soviet statistics}

For Estonia only the output side of national income is feasible in either period: consistent annual data on expenditure or factor incomes do not exist for the interwar republic, whose statistical office never compiled national accounts, and for the Soviet period the monetary aggregates that an expenditure or income estimate would require are precisely the figures that cannot be trusted. The series presented here are therefore constructed, in both periods, by extrapolating a benchmark estimate of GDP with constant-price indices of sectoral value-added, each sector weighted by its share of GDP in the benchmark year:
\[
    \mathit{GDP}_{t} = \mathit{GDP}_{\mathrm{bench}} \times \sum_{i=1}^{n} w_{i}\, I_{it},
\]
where $w_i$ is sector $i$'s share of benchmark GDP and $I_{it}$ its volume index in year $t$. The procedure is standard where expenditure and income data are sparse, and is the one used for Ireland, the European periphery and, in the companion studies, Lithuania and Latvia.\footnote{\autocite[880--3]{kenny_irelandgdp}; \autocite[22]{broadberryinmorys}; \autocite{klimantas2023lithuania}; \autocite{Klimantas2024}.}

For the Soviet period this approach also has an intellectual pedigree that bears directly on the reliability of the result. Soviet practice measured net material product at prevailing prices, excluded most services, double-counted intermediate transactions and, above all, understated inflation, so that official constant-price growth was too high: in Maddison's verdict, the official series `overstated growth, and understated inflation'.\footnote{\autocite[307]{Maddison1998b}.} The Western response, begun by Bergson in the 1950s and institutionalised by the CIA from the 1960s, was to rebuild Soviet GNP from the bottom up, repricing output at factor cost, imputing the missing services and, crucially, replacing the official value indices with indices built from physical output series weighted by value-added.\footnote{\autocite{bergson1950}; \autocite{bergson1961}; \autocite[48]{cia1982a}; \autocite[328]{Maddison1995a}.} The CIA estimates were attacked from both sides. Navrozov and others argued that by accepting Soviet physical data they still overstated growth; Khanin, working from Soviet materials in the 1980s, obtained much lower growth rates for the Stalin period but, as Harrison showed, results for 1950--87 that coincided with the CIA's at 3.8 per cent a year.\footnote{\autocite{assessingCIA1988}; \autocite[144, 159]{harrison1993}.} The CIA's own indices were reworked at 1982 prices for the Joint Economic Committee in 1990, the form in which they entered the later literature and the MPD.\footnote{\autocite{cia1990}.} The assessments by Becker and by Schroeder, the review by Maddison and the recent survey by Vonyó and Markevich all reach the same conclusion: the CIA series contain errors in both directions, most of the criticism `was not well founded', and the physical-output method remains the best available basis for measuring Soviet growth.\footnote{\autocite{becker1994}; \autocite{schroeder1995}; \autocite[309]{Maddison1998b}; \autocite[279--80]{vonyoinmorys2021}.} Ponomarenko's archival reconstruction of Russian national accounts, the fullest attempt so far to recast Soviet data in the System of National Accounts (SNA), reaches the same judgement from inside the Soviet statistical system. The physical output data `are apparently of high quality and can be used', whereas the official value indices were inflated by the pricing of new products; Khanin's alternative, in Ponomarenko's view, rests on a `rather primitive methodological base' compared with the CIA's.\footnote{\autocite[16--18, 73--4]{ponomarenko2002}.} The literature on the Soviet economy has grown rapidly since the opening of the archives, as Zhuravskaya, Guriev and Markevich document, but its quantitative core for the post-war decades still rests on physical output series of this kind.\footnote{\autocite[47--114]{zhuravskaya2024}; \autocite{voskoboynikov2021}.}

It matters, therefore, to be precise about which Soviet data this paper uses and which it avoids. The distortions documented in the literature were concentrated in value aggregates: in the national income and gross output indices at `comparable prices', in the reporting of hidden price increases as real growth, and in the double counting of the material product system.\footnote{\autocite[307]{Maddison1998b}; \autocite{vonyo2017}.} Khanin's classification of the official data, as Harrison summarises it, is that `the ease of distortion remained greater for value series than for physical volumes', and that distortion was concentrated in heterogeneous products, above all in machinery and construction, while physical series for homogeneous goods such as grain, milk, electricity, cement or railway freight are comparatively reliable.\footnote{\autocite[144, 152]{harrison1993}.} The sectoral indices below rest almost entirely on physical quantities of this kind, on employment and on inputs; the one place where value data at comparable prices enter, trade turnover, which also feeds the productivity adjustment for finance, is flagged and its bias bounded in Section~\ref{sec:limitations}. The claim, then, is not that Soviet statistics were honest, but that the particular statistics used here were recorded for operational planning and had little scope for manipulation.

\subsection{The benchmarks}

\textit{The 1938 benchmark.} Valge's estimate of current-price value-added for 1938 in agriculture, forestry and fishing (250.1 million kroon), industry (132.1 million), construction (14.0 million) and housing (29.0 million) is accepted without revision. His treatment of the remaining services is not. For all other services Valge assumed uniform value-added per worker and scaled a single figure by the service labour force, obtaining 122.9 million kroon in all, less than the contribution of industry and implausibly low for an economy at Estonia's level of urbanisation.\footnote{\autocite[2718]{valge2003all}.} This paper rebuilds service value-added from its income components, as in the Lithuanian and Latvian benchmark studies: compensation of employees is estimated from the labour force and average earnings of each service branch, and the remaining components, operating surplus, depreciation and net production taxes, are added in the proportions observed in the Latvian benchmark of 1935, the closest comparator with a full set of accounts.\footnote{\autocite{klimantas2019}; \autocite{norkuslatvia1935}.} Commerce yields 91.5 million kroon, transport and communications 32.2 million, public services 99.3 million, of which 21.4 million is defence, and household services 5.5 million: 228.5 million in all. Gross value-added rises to 654 million kroon and GDP at market prices, with 17.2 million of taxes less subsidies, to 671 million, or about 592 kroon per head against Valge's 502. The sectoral composition is reported in Table~\ref{tab:benchmarks}.

% tab:benchmarks
\begin{table}[htbp]
\centering\small
\caption[Benchmark GDP by sector: Estonia 1938 and 1990, current prices]{\textsc{Benchmark GDP by sector: Estonia 1938 and 1990, current prices}}
\label{tab:benchmarks}
\begin{threeparttable}
\begin{tabular}{lrrrr}
\toprule
& \multicolumn{2}{c}{1938} & \multicolumn{2}{c}{1990} \\
\cmidrule(lr){2-3}\cmidrule(lr){4-5}
Sector & m.\,\eek{} & \% & m.\,roubles & \% \\
\midrule
Agriculture, forestry \& fishing & 250.1 & 38.3 & 1,250.0 & 16.9 \\
Industry & 132.1 & 20.2 & 2,203.4 & 29.8 \\
Construction & 14.0 & 2.1 & 818.2 & 11.1 \\
Transport \& communications & 32.2 & 4.9 & 535.7 & 7.3 \\
Trade, finance, housing \& other commercial services & 120.5 & 18.4 & 1,322.3 & 17.9 \\
Public administration, defence, education \& health & 99.3 & 15.2 & 1,064.4 & 14.4 \\
Other services & 5.5 & 0.8 & 190.1 & 2.6 \\
\midrule
Gross value-added & 653.7 & 100.0 & 7,384.1 & 100.0 \\
Taxes less subsidies on products & 17.2 & & --- & \\
GDP at market prices & 670.9 & & 7,384.1 & \\
GDP per capita, national currency & 592.4 & & 4,705.7 & \\
\textit{GDP per capita, \intlnineteen{}} & \textit{2,685} & & \textit{8,425} & \\
\bottomrule
\end{tabular}
\begin{tablenotes}[para,flushleft]
\footnotesize
\textit{Sources:} 1938: Valge's estimates for agriculture, industry, construction and housing, with services re-estimated as described in Section~\ref{sec:method} (see online appendix). 1990: \cite[184--6]{bolotin1992}, converted to roubles with Bolotin's sectoral purchasing-power parities; services other than trade and transport split using the 1987 Soviet input--output table (\cite{wbinputoutput1987}). \textit{Notes:} the 1990 figure includes taxes less subsidies within sectoral values (`---': included in sectoral values). Population: mid-year.
\end{tablenotes}
\end{threeparttable}
\end{table}


To express this figure in international dollars, a bilateral comparison of consumer prices with Sweden was carried out for 1938, using prices of some fifty consumer goods and services in both countries, weighted by the structure of household expenditure in each. The Fisher price ratio is 1.79 Swedish kronor per Estonian kroon, against an official exchange rate of 1.06: the kroon bought in Estonia roughly 70 per cent more than its exchange value would have bought in Sweden. Applied to Swedish GDP per head of 1,931 kronor, the ratio implies that Estonian GDP per head was 54.8 per cent of Sweden's, and since Sweden stood at 4,901~\intlnineteen{} in 1938, Estonia's 1938 benchmark is \ee{1938}~\intlnineteen{}, or \eeb{1938}~\intltwenty{}.\footnote{\autocite{mpd2023}, with \intltwenty{} converted at 1.594 (\autocite[12]{bolt_zanden_newgdp2021}). The price comparison follows the procedure of \autocite{klimantas2019} and \autocite{norkuslatvia1935}; the underlying prices and weights are given in the online appendix.} The result is 2 per cent below Valge's 2,745, but the proximity is accidental: his domestic-currency estimate was too low by 15 per cent and his Finnish bridge too generous by a similar margin. Expressed as a share of Finland, the new benchmark gives \pctfin{1938} per cent for 1938 against Clark's 90 per cent; Norkus and Markevi\v{c}i\={u}t\.{e}'s independent approximation gives 71 per cent.\footnote{\autocite[154]{norkus_restorations}.}

\textit{The 1990 benchmark.} The Soviet-period series is anchored to Bolotin's estimate of ESSR GDP in 1990, the first systematic attempt to express the accounts of all fifteen republics in international dollars.\footnote{\autocite[181--92]{bolotin1992}.} Bolotin started from the republican GDP totals for 1990 compiled by Goskomstat in 1990--1, 1990 being the first year for which republican accounts on the SNA definition were compiled, and converted them into dollars with purchasing-power parities estimated, as he put it, `in practically full conformity' with the International Comparison Project (ICP). The parities, by sector of origin, were 0.40 dollars per rouble in agriculture, 1.80 in industry, 1.10 in construction, 2.80 in transport, 1.45 in trade and 3.30 in other services, or 1.82 for GDP as a whole.\footnote{\autocite[181, 184]{bolotin1992}.} For Estonia the result is \ee{bench1990}~\intlnineteen{} per head, the highest of the Baltic republics. Table~\ref{tab:benchmarks} shows the sectoral composition in roubles, with services other than trade and transport allocated by the 1987 Soviet input--output table;\footnote{\autocite{wbinputoutput1987}.} these rouble-price shares, domestic-price weights like the kroon shares of 1938, are the weights with which the volume indices are combined, the dollar conversion applying to the total alone.

Three features of this benchmark need to be stated plainly, because the reconstruction stands or falls with it. First, its level is conservative by the standards of the time. Bolotin's figure for the Soviet Union as a whole, 6,275 dollars per head, is below Maddison's own estimate for 1987, and the 1990 comparison on which it draws, a bilateral Austria--USSR price comparison conducted in the ICP framework, placed the USSR at 32 per cent of the United States where the CIA had placed it at 42 per cent.\footnote{\autocite[318--19]{Maddison1998b}.} Second, the parities are Union-wide, so that the conversion assumes that relative prices in Estonia were those of the Soviet Union as a whole. For the material sectors this is a reasonable approximation in a planned economy with administered prices; for services, where Bolotin's parity of 3.30 dollars per rouble reflects the very low rouble wages of the non-market sector, it is less certain. A separate question is the allocation of value-added among services with the 1987 input--output table: the resulting shares of trade and services in Estonian GDP at rouble prices (Table~\ref{tab:benchmarks}) are higher than those implied by the World Bank's 1990 sectoral data for Estonia, 16 per cent for agriculture, 47 per cent for industry and construction and 38 per cent for services.\footnote{\autocite{wbstatisticalhandbook1994}; \autocite[1--8]{wb1993estonia}.} The sensitivity of the series to this allocation is examined in Section~\ref{sec:limitations}. Third, there is no alternative. Ponomarenko's retrospective accounts, the most thorough archival reconstruction of Soviet republican national accounts, cover Russia alone, make no dollar conversion and, in their author's words, show that the official exchange rate `was greatly overvalued', so that any republican benchmark must rest on purchasing-power parities of Bolotin's kind.\footnote{\autocite[6, 15, 82]{ponomarenko2002}.} The CIA's 1982 benchmark cannot be decomposed by republic, and the 1987 input--output table understates services too severely to serve as a benchmark.\footnote{\autocite[66--89]{wbinputoutput1987}.} An independent Estonian benchmark with republic-specific prices is the natural next step; it would alter only the level to which the volume indices are anchored, not the indices themselves.

The choice of 1990 rather than an earlier year means that the Soviet-period series is extrapolated backwards for more than fifty years from end-of-period weights. The practical reason is that 1990 is the only year for which republican accounts on an SNA definition and purchasing-power parities exist together. Two further considerations favour it: the reporting of 1989--90, at the height of glasnost, was more transparent than that of any earlier Soviet year, the manipulation of earlier returns having been publicly acknowledged by Baltic statisticians in 1988;\footnote{\autocite{adlys1988}.} and a late benchmark captures the economy at its most diversified, so that branches that barely existed in 1940, electrical engineering, chemicals or oil-shale electricity, receive weights appropriate to their eventual size. The index-number consequences are treated in Section~\ref{sec:limitations}; the short answer is that they can be tested directly, because the Soviet-period series meets an independent series with 1938 weights in 1937--40.

\subsection{Volume indices and their sources}

Table~\ref{tab:sources} summarises the indicator and sources for each sector in each period; the online appendix documents every series year by year. Three general points precede the detail. First, the interwar indices are built from the publications of the Estonian State Statistical Central Bureau, the agricultural surveys of 1925, 1929 and 1939 and the compilations of Valge and of Klesment and Valge, and are weighted by the revised 1938 benchmark.\footnote{\autocite{valge2003all}; \autocite{klesment2007}.} Second, the Soviet-period indices are built from the statistical yearbooks of the ESSR, \textit{Eesti NSV rahvamajandus}, which were published from 1957 and give annual physical output from the mid-1950s. These are supplemented, for the years 1945--55 that the yearbooks cover incompletely or not at all, by the returns of the republican statistical administration held in the Estonian National Archives (fond R-10) and, for the weights within industrial branches, by an archival compilation of 1986 output values for several hundred products (fond R-1969).\footnote{See Table~\ref{tab:sources}. The archival sources are listed in the online appendix, section A.2.} Archival material is thus used systematically, not merely to fill gaps. Third, the series cover the civilian economy as reported to the republican statistical office. Military production, organised through all-Union enterprises and never reported at the republican level, is not captured: Estonia hosted fewer such plants than Latvia, but the largest, the Dvigatel works in Tallinn and the uranium plant at Sillam\"{a}e, lay outside republican statistics, an unquantifiable source of understatement;\footnote{\autocite[150]{kukk2005}; \autocite[66]{ahderajasalu1992}.} the size of the Soviet forces stationed in the republic is likewise unknown and is treated as a source of bias in Section~\ref{sec:limitations}.

% tab:sources
\begin{table}[htbp]
\centering\footnotesize
\caption[Volume indicators and sources by sector]{\textsc{Volume indicators and sources by sector}}
\label{tab:sources}
\begin{threeparttable}
\begin{tabular}{>{\RaggedRight}p{2.6cm}>{\RaggedRight}p{5.3cm}>{\RaggedRight}p{5.3cm}c}
\toprule
Sector & Interwar series, 1919--40 (1938 weights) & Soviet-period series, 1937--95 (1990 weights) & App. \\
\midrule
Crops & Physical output of 15 crops weighted by 1938 value-added (Valge; 1939 agricultural survey); yields from \cite{klesment2007} & Physical output; weights from World Bank (1996) and Eurostat agricultural accounts; 1945--65 minor crops interpolated between quinquennial observations & A.1 \\
Livestock & Physical output of meat, milk, eggs and wool; 1938 value-added weights & Physical output; during 1941--44 animal numbers multiplied by pre-war yields & A.1 \\
Forestry and fishing & Timber removals and fish catch; timber for 1919--21 and fish for 1919--27 extrapolated from the crops-and-livestock index & Timber removals and fish catch; 1941--44 interpolated & A.1 \\
Industry & Employment adjusted for real wages (no usable annual output data; Valge's current-price series cannot be deflated by branch) & 101 physical product series in 11 branches; branch weights from the 1990 benchmark, product weights from 1986 archival output data (ERA) & A.2 \\
Construction & Domestic cement use (production plus net imports), as in Valge; 1919--21 from population & Weighted index of material inputs (1987 Soviet input--output weights), following the CIA procedure & A.3 \\
Transport and communications & Railway passenger- and tonne-kilometres, merchant fleet tonnage, postal items, telegrams and telephone calls & Freight tonne-kilometres and passenger-kilometres for all modes; postal, telephone and broadcasting indicators; 1941--45 interpolated between 1940 and the first post-war observation, with 1943--45 movements from Lithuanian ratios (Section~\ref{sec:method}) & A.4 \\
Trade, finance and insurance & Marketed value-added of agriculture and industry plus deflated imports (`commerce' index); finance assumed to move with commerce & Soviet retail and wholesale turnover at comparable prices (archival and published); finance and insurance from employment with a productivity adjustment & A.5 \\
Housing & Valge's housing value-added for 1929--38, extrapolated with the urban population & Housing stock from construction and population data & A.5 \\
Public administration, defence, health and education & Civil services assumed to move with GDP; defence from military personnel (15,900 in 1921--39) & Employment by branch (Soviet yearbooks, archival records); defence from the 1940 and 1990 totals & A.6 \\
Other services & Total population & Employment & A.6 \\
\bottomrule
\end{tabular}
\begin{tablenotes}[para,flushleft]
\footnotesize
\textit{Sources:} Interwar: \cite{valge2003all}; \cite{klesment2007}; interwar statistical yearbooks and agricultural surveys. Soviet period: \textit{Eesti NSV rahvamajandus} statistical yearbooks (1957--89), \textit{Eesti statistika aastaraamat} (1990--98), Estonian National Archives (ERA, fond R-10 and R-1969), German occupation statistical bulletins (\textit{Statistische Berichte f\"ur das Ostland}, 1942--44), and World Bank (1993, 1996). Full source lists by series are given in the online appendix, sections A.1--A.6.
\end{tablenotes}
\end{threeparttable}
\end{table}


\textit{Agriculture, forestry and fishing.} In both periods the index combines physical output of crops (fifteen products, including fodder) and of livestock products (meat, milk, eggs and wool) with timber removals and the fish catch, each sub-index weighted by its benchmark-year value-added. Interwar output comes from Klesment and Valge's compilation and the agricultural surveys, Soviet-period output from the yearbooks and, for the early post-war years, the archives, where some minor crops are reported only quinquennially and have been interpolated.\footnote{\autocite{klesment2007}; online appendix, section A.1.} The weights for the Soviet period derive from the World Bank's 1990 agricultural accounts for Estonia and, for the split within grains, from the Eurostat economic accounts for agriculture.

\textit{Industry.} No annual volume index of interwar industrial output exists, and Valge's current-price series cannot be deflated by branch for want of producer price indices. The interwar industrial index is therefore built from industrial employment adjusted for the movement of real wages, a recognised second-best when the share of industry in GDP is modest.\footnote{\autocite[22]{broadberryinmorys}; \autocite[427--8]{thomasfeinstein2003}.} For the Soviet period, by contrast, physical output series could be assembled for 101 products in eleven branches: electricity, petroleum products, other fuels including oil shale, metals, chemicals, machinery and metalworking, wood and paper, building materials, light industry, food and a residual. Products are aggregated to branches with the 1986 archival weights and branches to industry with the weights of the 1990 benchmark, following the CIA's procedure for Soviet GNP.\footnote{\autocite[48]{cia1982a}; \autocite[312]{maddison2010}.}

\textit{Construction.} The interwar index follows Valge in tracking domestic consumption of cement, production plus net imports; it is volatile in the 1920s, when timber construction, which the index does not capture, was still important.\footnote{\autocite[Table 35]{valge2003all}.} The Soviet-period index is a weighted index of the material inputs into construction, with weights from the 1987 input--output table, which is the CIA's procedure.\footnote{\autocite[66--89]{wbinputoutput1987}.}

\textit{Transport and communications.} The interwar index combines railway passenger- and tonne-kilometres, the tonnage of the merchant fleet, and postal items, telegrams and telephone calls. The Soviet-period transport index combines freight tonne-kilometres and passenger-kilometres across all modes; communications combine postal and telegraph traffic, telephone calls and radio and television sets in use.

\textit{Trade, finance and housing.} For the interwar years a single `commerce' index tracks the marketed value-added of agriculture and industry together with deflated imports, and finance and insurance are assumed to move with it. Under planning, industrial products circulated by allocation rather than through markets, so the Soviet-period trade index is built instead from the retail and wholesale turnover statistics at comparable prices compiled by the republican statistical office, most of them archival; these are the only value series used in the reconstruction, and the hidden inflation they may contain is bounded in Section~\ref{sec:limitations}. Finance and insurance follow employment, adjusted for productivity with the ratio of trade turnover to trade employment, as the CIA did; housing follows the housing stock, derived from construction and population.\footnote{\autocite[319]{Maddison1998b}; \autocite[48]{cia1982a}.}

\textit{Public and other services.} Interwar public services other than defence are assumed to move with the rest of the economy, since no consistent series of civil employment exists; defence follows military personnel, which is known for 1919--20 and 1939 and held constant in the intervening years. Soviet-period public administration, health and education follow employment by branch, as in the CIA's accounts, with defence derived from the 1940 and 1990 totals. Other, mainly domestic, services follow population before 1940 and employment thereafter.

\textit{The war years.} Estonia was under German civil administration from 1941 to 1944, and the \textit{Reichskommissariat Ostland} published statistical bulletins that cover agriculture, the extractive industries and employment but little else. Where these fail, forestry, fishing, transport and several service branches are interpolated between 1940 and the first post-war observation, and the movement between 1943 and 1945, for which no Estonian figures could be found, is taken from the ratios observed in Lithuania, where documentation is better. The estimates for 1941--5 are therefore `rough' in Feinstein's sense, with a margin of 15--25 per cent, and the 1944 trough is an order of magnitude rather than a measurement.\footnote{\autocite[21]{feinstein1972}.} Because the Soviet-period series is extrapolated backwards through the war, an error in the wartime linkage would shift its level in 1937--40, so that the overlap tests the linkage together with the rest of the extrapolation; the spliced series, which follows the interwar series to 1940, is affected only through the 1940 average.

\subsection{Population}

Estonia's borders changed once, when the area east of Narva and most of Petseri county, together about 6 per cent of the interwar population, were transferred to the Russian republic in 1944--5. The interwar figures refer to the borders of the republic and the Soviet-period and modern figures to present-day borders, which is the convention of the MPD; the transferred districts were rural and poorer than the national average, so that the break in income per head is likely to be small, and no adjustment has been made.\footnote{\autocite[30]{bolt_zanden_newgdp2021}.} The population series rests on the censuses of 1922, 1934, 1959, 1970, 1979 and 1989 and the year-end estimates of Statistics Estonia and the `Baltic 100' project, with mid-year figures as the average of consecutive year-ends; wartime losses, deportations and the flight of some 70,000 people to the West in 1944 are taken from the demographic literature, and the Soviet garrison of 1940--1 is added to the resident population (Table~\ref{tab:series}).

\subsection{Aggregation and splicing}

Each of the two new series is first compiled in constant benchmark-year prices and then expressed in \intlnineteen{} through its benchmark. The interwar series covers 1919--40 and the Soviet-period series 1937--95, so that the two overlap for four years. They are spliced at 1940 by taking the average of the two estimates for that year, and the Soviet-period series is in turn spliced at 1991 to a modern series obtained by backcasting the MPD's 2011 benchmark with the World Bank's constant-price growth rates of GDP per capita, which are used in preference to the MPD's own 1990--2010 growth rates for reasons explained in Section~\ref{sec:limitations}.\footnote{\autocite{mpd2023}; \autocite{wdi2024}.} The 1913 figure is the benchmark of Norkus and Markevi\v{c}i\={u}t\.{e}, 2,058~\intlnineteen{}, taken as published; there are no estimates for 1914--18, and the change between 1913 and 1919 therefore combines their income-side approximation with the new production-side estimate and should be read with that caveat. Following Bolt and van Zanden, the series is compiled in \intlnineteen{} and converted to \intltwenty{} at 1.594, and both versions are reported.\footnote{\autocite[12, 30]{bolt_zanden_newgdp2021}.}
